\documentclass[11pt]{article}
\usepackage{color}

%----------Packages----------
\usepackage{amsmath}
\usepackage{amssymb}
\usepackage{amsthm}
\usepackage{dsfont}
\usepackage{mathrsfs}
\usepackage{stmaryrd}
\usepackage[all]{xy}
\usepackage[mathcal]{eucal}
\usepackage{verbatim}  %%includes comment environment
\usepackage{fullpage}  %%smaller margins
\usepackage{multicol}
\usepackage{hyperref}
\usepackage{graphicx}

%Packages added to get the code block
\usepackage{listings}
\usepackage{xcolor}
\usepackage{tikz-cd}

\definecolor{codegreen}{rgb}{0,0.6,0}
\definecolor{codegray}{rgb}{0.5,0.5,0.5}
\definecolor{codepurple}{rgb}{0.58,0,0.82}
\definecolor{backcolour}{rgb}{0.95,0.95,0.92}

\lstdefinestyle{mystyle}{
    backgroundcolor=\color{backcolour},   
    commentstyle=\color{codegreen},
    keywordstyle=\color{magenta},
    numberstyle=\tiny\color{codegray},
    stringstyle=\color{codepurple},
    basicstyle=\ttfamily\footnotesize,
    breakatwhitespace=false,         
    breaklines=true,                 
    captionpos=b,                    
    keepspaces=true,                 
    numbers=left,                    
    numbersep=5pt,                  
    showspaces=false,                
    showstringspaces=false,
    showtabs=false,                  
    tabsize=2
}

\lstset{style=mystyle}


%---------Commands----------
\def\ZZ{\mathbb{Z}}
\def\QQ{\mathbb{Q}}
\def\PP{\mathbb{P}}
\def\NN{\mathbb{N}}
\def\OO{\mathcal{O}}
\def\AA{\mathbb{A}}
\def\BB{\mathbb{B}}
\def\FF{\mathbb{F}}
\def\LL{\mathbb{L}}
\def\RR{\mathbb{R}}
\def\CC{\mathbb{C}}
\def\GG{\mathbb{G}}
\def\lcm{\mathrm{lcm}}
\def\ord{\mathrm{ord}}
\def\rep{\mathrm{Rep}}
\def\dom{\mathrm{Dom}}
\def\ve{\varepsilon}
\def\gal{\mathrm{Gal}}
\def\emb{\mathrm{Emb}}

\newcommand{\dlegendre}[2]{\ensuremath{\left( \frac{#1}{#2} \right) }}
\newcommand{\tlegendre}[2]{\ensuremath{\big( \frac{#1}{#2} \big) }}

%----------Theorems----------
\theoremstyle{definition}
\newtheorem{theorem}{Theorem}[section]
\newtheorem{proposition}[theorem]{Proposition}
\newtheorem{lemma}[theorem]{Lemma}
\newtheorem{corollary}[theorem]{Corollary}
\newtheorem{definition}[theorem]{Definition}
\newtheorem{example}[theorem]{Example}
\newtheorem*{definition*}{Definition}
\newtheorem{nondefinition}[theorem]{Non-Definition}
\newtheorem{problem}{Problem}[section]
\newtheorem*{tip}{Pro Tip}
\newtheorem*{remark}{Remark}
\numberwithin{equation}{subsection}

\newenvironment{solution}
{%
  \vspace{0.5em}%
  \hrule%
  \vspace{0.2em}%
  \noindent\textbf{Solution:}\!\!\!\!\!
  \pushQED{\qed}
}
{%
  \popQED
  \vspace{.2em}
  \hrule%
  \vspace{0.75em}%
}

\newcommand{\separate}{
    \vspace{.2in}
    \hrule
    \vspace{.2in}
}


\newcommand{\head}[1]{
	\begin{center}
		{\large #1}
		\vspace{.05 in}
	\end{center}
	\bigskip 
}

\newcommand{\blue}[1]{\textcolor{blue}{#1}}
\newcommand{\red}[1]{\textcolor{red}{#1}}

%Problem Set Data
%number = 8
%course = 261
%semester = Fall 2025
%path = /Users/thyde/Documents/Teaching/Vassar/Semesters/2025 Fall/Math 261 Fall 2025/Problem Sets/Problem Set 8

\begin{document}
\pagestyle{plain}

%%---  sheet number for theorem counter
\setcounter{section}{2}   

\head{MATH 999, Summer 2026\\ Note 2: The Extension Theorem\\
Trevor Hyde}

%HEADER

\begin{center}
   \emph{It is with surprise and wonderment that over the years I discovered (or rather, doubtless, rediscovered) the prodigious, truly inexhaustible richness, the unsuspected depth of this theme, apparently so anodyne. 
   I believe I feel a central sensitive point there, a privileged point of convergence of the principal currents of mathematical ideas, and also of the principal structures and visions of things which they express...}\\
   \vspace{.1in}
   
   Alexander Grothendieck on the theme of Galois theory in mathematics, from \emph{Esquisse d'un Programme}.
\end{center}
\vspace{.1in}

In the previous note we discussed algebraic extensions of a field $K$, including the ultimate algebraic extension $\overline{K}/K$, which contains all the roots of every polynomial $f(x) \in K[x]$.
We often do not need to be so greedy and instead only want to consider extensions which contain all the roots of some given polynomial.
That leads us to the notion of a \emph{splitting field}.

\begin{definition}
    We say $L/K$ is a \textbf{splitting extension} of a monic polynomial $f(x) \in K[x]$ if there are $\alpha_1,\ldots, \alpha_n \in L$ such that
    \[
        f(x) = (x - \alpha_1)\cdots (x - \alpha_n),
    \]
    and $L = K(\alpha_1,\ldots, \alpha_n)$.
    In other words, $f(x)$ splits completely in $L$ and $L$ is the smallest extension of $K$ for which this occurs.
\end{definition}

For example, $\CC/\RR$ is a splitting extension of $x^2 +1 \in \RR[x]$ since $x^2 + 1 = (x - i)(x + i)$ and $\CC = \RR(i, -i)$.
Note that in this case we did not actually need both of the roots of $x^2 + 1$ to generate $\CC$, but that's okay: the point is just that $x^2 + 1$ cannot split completely over any smaller extension.

Splitting extensions are sensitive to both fields: we just saw that $\CC/\RR$ was the splitting extension for $x^2 + 1$ viewed as a polynomial in $\RR[x]$, however $\CC/\QQ$ is not a splitting extension for $x^2 + 1 \in \QQ[x]$.
The reason is that $\QQ(i,-i) = \QQ(i)$ is \emph{much} smaller than $\CC$; the field $\QQ(i)$ is countable.
\footnote{
If $L/K$ is a splitting extension, we often call $L$ a \textbf{splitting field}, but just beware that when doing so the field $K$ is implicit in that. 
}

Splitting extensions always exist.
If $f(x) \in K[x]$, then all the roots $\alpha_1,\ldots, \alpha_n$ of $f(x)$ live in $\overline{K}$ and thus $K \subseteq K(\alpha_1,\ldots, \alpha_n) \subseteq \overline{K}$ gives us a splitting extension of $K$.
But here's a philosophical question:
If we construct splitting extensions of $f(x) \in K[x]$ in two different ways $L_1/K$ and $L_2/K$, in what sense are these two fields ``the same''?
They don't need to be on the nose equal.
For example, consider $x^2 - 2 \in \QQ[x]$.
One way to construct a splitting field is to first construct the real numbers \emph{a la} analysis and then to prove that 2 has a square root $\sqrt{2} \in \RR$.
Then we have $\QQ \subseteq \QQ(\sqrt{2}) \subseteq \RR$ and $\QQ(\sqrt{2})/\QQ$ is a splitting extension of $x^2 - 2$.
On the other hand, we can formally construct a splitting field via $L := \QQ[y]/\langle y^2 - 2\rangle$.
We have
\[
    x^2 - 2 \equiv (x - y)(x + y) \bmod y^2 - 2.
\]
The field $L$ is not literally contained in the real numbers, so $L \neq \QQ(\sqrt{2})$.
However, as we saw in a previous note, these two fields are isomorphic, with an isomorphism induced by the evaluation map.

This example highlights what is true in general: any two splitting extensions of $f(x) \in K[x]$ are isomorphic.
Proving this takes some effort.
We will actually prove a stronger, more abstract result called the Extension Theorem\footnote{This theorem does not have an agreed upon name, but I think Extension Theorem is suitable.} from which the uniqueness of splitting fields may be deduced.
The Extension Theorem is a fundamental technical building block for Galois theory.
Before we get to it, a lemma and a couple definitions.

\begin{definition}
    Suppose $\phi : X_1 \to X_2$ is a function between sets $X_1$ and $X_2$.
    Suppose that $X_1 \subseteq Y_1$ and $X_2 \subseteq Y_2$ are larger sets and that $\psi : Y_1 \to Y_2$.
    We say that $\psi$ \textbf{extends} $\phi$ if $\psi(x) = \phi(x)$ for all $x \in X_1$.
\end{definition}

\begin{lemma}
\label{lemma: extend}
    Let $\phi : K_1 \to K_2$ be an isomorphism of fields.
    Let $f_1(x) \in K_1[x]$ be an irreducible polynomial and let $f_2(x) \in K_2[x]$ be the polynomial we get by applying $\phi$ to each coefficient of $f_1$.
    Let $\alpha_1$ be a root of $f_1$ in some extension of $K_1$ and let $\alpha_2$ be a root of $f_2$ in some extension of $K_2$.
    Then there is a unique isomorphism $\psi : K_1(\alpha_1) \to K_2(\alpha_2)$ such that $\psi$ extends $\phi$ and $\psi(\alpha_1) = \alpha_2$.
\end{lemma}

\begin{proof}
    \blue{First prove that $f_2(x)$ is irreducible in $K_2[x]$.}
    Let $\phi' : K_1[x] \to K_2[x]$ be the isomorphism of polynomial rings where $\phi'(g(x))$ is the polynomial we get by applying $\phi$ to all the coefficients of $g(x)$ (\blue{check that this really is ring isomorphism}).
    Then $\phi'(\langle f_1(x)\rangle) = \langle f_2(x)\rangle$, hence $\phi'$ induces an isomorphism of fields
    \[
        \overline{\phi'} : \frac{K_1[x]}{\langle f_1(x)\rangle} \xrightarrow{\sim} \frac{K_2[x]}{\langle f_2(x)\rangle}.
    \]
    Furthermore we have the isomorphisms induced by evaluation maps
    \[
        \overline{\ve}_{\alpha_i} : \frac{K_i[x]}{\langle f_i(x)\rangle} \xrightarrow{\sim} K_i(\alpha_i).
    \]
    Now let $\psi := \overline{\ve}_{\alpha_2} \overline{\phi'}\overline{\ve}_{\alpha_1}^{-1}$, where we denote composition of maps by juxtaposition.
    To diagram below helps us ``visualize'' the map we have just constructed.
    \footnote{This is an example of a \emph{commutative diagram}, which expresses an identity of map compositions.
    The idea is that you can follow any two paths in the diagram between the same starting and ending point and the result will be the same.
    If you move along an arrow in the opposite direction in which it is pointing, you take the inverse of that map.
    This diagram captures the definition of $\psi$, the bottom arrow, by taking the alternate path around the top of the diagram.
    These sorts of commutative diagrams show up all the time in algebra.}
    \[
        \begin{tikzcd}[row sep=2.5em, column sep=4em]
        \dfrac{K_1[x]}{\langle f_1(x) \rangle} \arrow[r, "\overline{\phi'}"] \arrow[d, "\overline{\ve}_{\alpha_1}"'] & \dfrac{K_2[x]}{\langle f_2(x) \rangle} \arrow[d, "\overline{\ve}_{\alpha_2}"] \\
        K_1(\alpha_1) \arrow[r, "\psi"'] & K_2(\alpha_2)
        \end{tikzcd}
    \]
    Let's check that the map $\psi$ has the two properties we want.
    \blue{First show that $\psi$ extends $\phi$} and then \blue{show that $\psi(\alpha_1) = \alpha_2$.}
    Don't overthink it, just carefully unpack the definitions of each map that makes up the composition defining $\psi$.
    That finishes the proof of existence.

    And for uniqueness, suppose that $\psi_1, \psi_2$ both extend $\phi$ and both send $\alpha_1$ to $\alpha_2$.
    An element $\beta \in K_1(\alpha_1)$ is, by definition $g_1(\alpha_1)$ where $g_1(x) = a_0 + a_1x + \ldots + a_d x^d \in K_1[x]$.
    Then for $i = 1, 2$ we have
    \begin{align*}
        \psi_i(\beta) &= \psi_i(a_0 + a_1\alpha_1 + \ldots + a_d\alpha_1^d)\\
        &= \psi_i(a_0) + \psi_i(a_1)\psi_i(\alpha_1) + \ldots + \psi_i(a_d)\psi_i(\alpha_1)^d\\
        &= \phi(a_0) + \phi(a_1)\alpha_2 + \ldots + \phi(a_d)\alpha_2^d.
    \end{align*}
    Hence $\psi_1(\beta) = \psi_2(\beta)$.
    This implies that $\psi_1 = \psi_2$.
\end{proof}

\begin{definition}
    A (monic)\footnote{There are a few definitions where I snuck in a monic assumption.
    This is just for a minor notational simplicity.
    The definitions also make sense for general polynomials, but you'll need to make minor and hopefully obvious tweaks.} polynomial $f(x) \in K[x]$ is \textbf{separable} if all the roots of $f(x)$ in $\overline{K}$ are distinct.
    That is
    \[
        f(x) = (x - \alpha_1)\cdots (x - \alpha_d)
    \]
    for some $\alpha_i \in \overline{K}$ where $\alpha_i \neq \alpha_j$ when $i \neq j$.
\end{definition}

Most of the time, irreducible polynomials are separable.
However, there are some weird situations in characteristic $p$ where this is not the case.
We will discuss this later.
For now just remember that separable means that we can separate all the roots from one another.

\begin{theorem}[Extension Theorem]
\label{thm: extension theorem}
    Let $\phi : K_1 \to K_2$ be an isomorphism of fields.
    Let $f_1(x) \in K_1[x]$ be a polynomial and let $f_2(x) \in K_2[x]$ be the polynomial we get by applying $\phi$ to each coefficient of $f_1$.
    Let $L_1/K_1$ be a splitting extension of $f_1$ and let $L_2/K_2$ be a splitting extension of $f_2$.
    Then
    \begin{enumerate}
        \item (Existence) There exists an isomorphism $\psi : L_1 \to L_2$ which extends $\phi : K_1 \to K_2$.

        \item (Multiplicity) The total number of distinct isomorphisms $\psi : L_1 \to L_2$ extending $\phi$ is at most $[L_1 : K_1]$.
        If $f(x)$ is separable, then the number of extensions is exactly $[L_1 : K_1]$.
    \end{enumerate}
\end{theorem}

\begin{proof}
    We will prove this by induction on $d := [L_1 : K_1]$.
    Let's start with the base case, $d = 1$.
    If $d = 1$, then $L_1 = K_1$.
    So we can simply let $\psi = \phi$.
    There is no other way to extend $\phi$ so the extension is unique.

    Now suppose that $d > 1$ and that the result is true for all splitting extensions of lower degree.
    Since $d > 1$, the polynomial $f_1(x)$ does not split completely over $K_1$, hence must have some irreducible factor $g_1(x)$.
    Let $g_2(x) \in K_2[x]$ be the polynomial we get by applying $\phi$ to each of the coefficients of $g_1(x)$.
    Note that $g_2(x)$ is irreducible in $K_2[x]$ and is a factor of $f_2(x)$ (\blue{but don't take my word for it, think it through}).
    Let's say $e := \deg g_1 = \deg g_2$ and let $\alpha_1 \in L_1$ be a root of $g_1$.
    Since $f_2$ splits completely over $L_2$, so does $g_2$.
    Let $\alpha_{2,1}, \alpha_{2,2},\ldots, \alpha_{2,e'}$ be all the distinct roots of $g_2$ in $L_2$.
    We always have $e' \leq e$ (\blue{why?}).
    If $f_1$ and hence $f_2$ are separable, then so are their factors $g_1$ and $g_2$.
    In that case $e' = e$.
    For each $1 \leq j \leq e'$, Lemma \ref{lemma: extend} implies that $\phi$ has a unique extension $\phi_j : K_1(\alpha_1) \to K_2(\alpha_{2,j})$ such that $\phi_j(\alpha_1) = \alpha_{2,j}$.
    Furthermore, if $\phi' : K_1(\alpha_1) \to L_2$ is any extension of $\phi$, then applying $\phi'$ to $g_1(\alpha_1) = 0$ gives us $g_2(\phi'(\alpha_1)) = 0$.
    Hence $\phi'$ must send $\alpha_1$ to a root of $g_2$.
    Thus the $\phi_j$ are the only possible extensions of $\phi$.

    Now if we fix some $j$, let's suppose that we consider $K_1' := K_1(\alpha_1)$ and $K_2' := K_2(\alpha_{2,j})$ as our new base fields and $\phi_j : K_1' \to K_2'$ our now starting isomorphism.
    Then $f_1(x) \in K_1'[x]$ and since $\phi_j$ extends $\phi$, we still have that $f_2(x) \in K_2'[x]$ is the result of applying $\phi_j$ to the coefficients of $f_1$.
    Then $L_1/K_1'$ is still a splitting extension of $f_1$ and so is $L_2/K_2'$ a splitting extension of $f_2$, but the tower law implies that $[L_1 : K_1'] < [L_1 : K_1'][K_1' : K_1] = [L_1 : K_1] = d$.
    Therefore the induction hypothesis kicks in to tell us that $\phi_j$ has at least one and at most $[L_1 : K_1']$ extensions to an isomorphism $\psi_j : L_1 \to L_2$.
    If $f_1$ is separable, the induction hypothesis implies that $\phi_j$ has exactly $[L_1 : K_1']$ extensions.
    Note that since $\psi_j$ extends $\phi_j$, which itself extends $\phi$, we have that each $\psi_j$ extends $\phi$.
    Since this is true for each $\phi_j$ and there are $e' \leq [K_1' : K_1]$ possible maps $\phi_j$, that gives us a total of at most $e'[L_1 : K_1'] \leq [L_1 : K_1'][K_1' : K_1] = [L_1 : K_1]$ extensions of $\phi$ and in the separable case, precisely this many.
    That completes our induction.
\end{proof}

This proof is our most challenging yet, it may take a few reads before it sinks in.
Here is our first application.
If $K$ is a field and $L_1/K$ and $L_2/K$ are field extensions, then we say a homomorphism $\phi : L_1 \to L_2$ is a \textbf{$K$-homomorphism} if $\phi(a) = a$ for all $a \in K$.
Equivalently, $\phi$ extends the identity map on $K$.

\begin{corollary}[Uniqueness of Splitting Extensions]
\label{cor: uniqueness of splitting extensions}
    Let $f(x) \in K[x]$ be a polynomial.
    If $L_1/K$ and $L_2/K$ are two splitting fields of $f$, then there exists at least one and at most $[L_1 : K]$ distinct $K$-isomorphisms $\psi : L_1 \to L_2$.
    If $f$ is separable, there are exactly $[L_1 : K]$ distinct $K$-isomorphisms between $\psi : L_1 \to L_2$.
\end{corollary}

\begin{proof}
    We directly apply the Extension Theorem to the identity map $\phi : K \to K$.
\end{proof}

Even though Corollary \ref{cor: uniqueness of splitting extensions} was our main goal, in order to prove it via induction we needed to state the result in a more general setting as in the statement of Theorem \ref{thm: extension theorem}.
Since the conclusion is more general, we might as well make that the theorem.

Just a reminder of the philosophy of algebra: if two objects are isomorphic, we treat them as essentially the same.
Hence Corollary \ref{cor: uniqueness of splitting extensions} is telling us that any two ways you might construct a splitting field of a polynomial $f(x) \in K[x]$, those splitting fields are always $K$-isomorphic.
The fact that it's a $K$-isomorphism (as opposed to just a plain old isomorphism) is important here: this means that the isomorphism $\psi$ between the splitting fields fixes everything in $K$, so really the two splitting fields look identical from the point of view of $K$.
If it were not a $K$-isomorphism, then $\psi$ could potentially twist around $K$, replacing $f(x)$ with a totally different polynomial.
In light of this result, we may simply talk about \emph{the} splitting extension of $f(x) \in K[x]$ and let this mean any splitting extension without worrying about how it was constructed; they're all essentially the same.

Corollary \ref{cor: uniqueness of splitting extensions} also tells us something about how many $K$-isomorphisms there are between splitting fields.
To see the importance of this, we'll need a corollary of the corollary.
First, an important definition.
Recall that an automorphism of an object $L$ is an isomorphism $\psi : L \to L$ from $L$ to itself.

\begin{definition}
    Let $L/K$ be a finite extension.
    The \textbf{Galois group} of $L/K$ is the group of all $K$-automorphisms $\psi : L \to L$ with composition as the group operation.
    We let $\gal(L/K)$ denote the Galois group.
\end{definition}

\begin{problem}
    \blue{Verify that $\gal(L/K)$ satisfies the group axioms.}
\end{problem}

The Galois group of an extension should be thought of as the group of symmetries of the extension.
It consists of all ways $L$ can be transformed which cannot be distinguished from the point of view of $K$.
The following result bounds the size of the Galois group of a splitting extension, and even better when our polynomial is separable.

\begin{corollary}
\label{cor: Galois order}
    Let $f(x) \in K[x]$ be a polynomial and let $L/K$ be a splitting extension.
    Then $|\gal(L/K)| \leq [L : K]$ with equality holding whenever $f(x)$ is separable.
\end{corollary}

\begin{proof}
    This is a direct application of Corollary \ref{cor: uniqueness of splitting extensions} with $L_1 = L_2$.
\end{proof}

At this point, Galois groups are quite mysterious.
We have given a formal, abstract definition of them without yet seeing any examples.
The Extension Theorem (via Corollary \ref{cor: Galois order}) tells us that for splitting extensions, the Galois group is always finite.
Actually, the upper bound we got in Corollary \ref{cor: Galois order} holds for all finite extensions.

\begin{problem}
    Our goal in this problem is to prove that if $L/K$ is a finite extension, then $|\gal(L/K)| \leq [L : K]$.
    In order to get this we will need to prove something more general, similar to how we proved the Extension Theorem in order to get Corollary \ref{cor: uniqueness of splitting extensions} and Corollary \ref{cor: Galois order}.
    Given two extensions $L_1/K$ and $L_2/K$, let $\emb_K(L_1,L_2)$ denote the set of all injective $K$-homomorphisms (also called \textbf{$K$-embeddings}).
    We will first prove that $|\emb_K(L,\overline{K})| \leq [L : K]$ for any finite extension $L/K$, and then achieve our goal as a corollary.
    \begin{enumerate}
        \item \blue{Prove that any finite extension $L/K$ may be expressed as $L = K(\alpha_1,\ldots, \alpha_n)$ for finitely many elements $\alpha_i \in L$.}

        \item We now use induction on $n$ to prove that any finite extension $L/K$ where $L = K(\alpha_1,\ldots, \alpha_n)$ we have $|\emb_K(L,\overline{K})| \leq [L : K]$.
        \begin{enumerate}
            \item \blue{Prove the base case.}
            \textbf{Hint:}
            First show that any $K$-embedding of $K(\alpha_1)$ into $\overline{K}$ is determined by where it sends $\alpha_1$.
            What are the options for where to send $\alpha_1$?
            This has something to do with the minimal polynomial $m_1(x) \in K[x]$ of $\alpha_1$.

            \item \blue{Prove the inductive step.}
            \textbf{Hint:}
            Try to use the inductive hypothesis to reduce this step back to a similar argument as the base case.
        \end{enumerate}

        \item Any finite extension $L/K$ is contained in some algebraic closure $K \subseteq L \subseteq \overline{K}$.
        \blue{Prove that $\gal(L/K) \subseteq \emb_K(L,\overline{K})$ and use this to prove that $|\gal(L/K)| \leq [L : K]$.}
    \end{enumerate}
\end{problem}

Thus we know that Galois groups of finite extensions are always finite and bounded in size by the degree of the extension.
Corollary \ref{cor: Galois order} tells us that splitting fields of separable polynomials have the largest possible Galois groups.
In other words, these are the \emph{most symmetric} extensions.
We give these extensions a special name.

\begin{definition}
    A \textbf{Galois extension} $L/K$ is a splitting extension of a separable polynomial $f(x) \in K[x]$.
\end{definition}

The following is a direct consequence of Corollary \ref{cor: Galois order}.

\begin{corollary}
    If $L/K$ is a Galois extension, then $|\gal(L/K)| = [L : K]$.
\end{corollary}

Let's consider some examples.

\begin{problem}\,
    \begin{enumerate}
        \item \blue{Let $K$ be a field of characteristic not equal to 2.
        Prove that any quadratic extension $L/K$ is Galois.}

        \item \blue{Now suppose that $K$ has characteristic 2.
        In the previous note we classified the quadratic extensions of $K$.
        Which of these quadratic extensions is Galois?}

        \item Let $\sqrt[3]{2}$ be a real cube root of 2.
        \blue{Prove that $\QQ(\sqrt[3]{2})/\QQ$ is not a Galois extension.}
    \end{enumerate}
\end{problem}

\begin{problem}
    Although the theoretical development of Galois theory is quite clean and streamlined, when we turn to examples things get gnarly.
    We must use every trick we can think of to squeeze the truth out of polynomials.
    For this problem, let $f(x) = x^3 + x^2 - 2x - 1 \in \QQ[x]$.

    \begin{enumerate}
        \item \blue{Prove that a cubic polynomial $g(x) \in K[x]$ is irreducible if and only if $g(x)$ has no roots in $K$.}
        
        \item \blue{Prove that $f(x)$ has no roots in $\QQ$ and use this to conclude that $f(x)$ is irreducible in $\QQ[x]$.}

        \item \blue{Prove that if $\alpha$ is a root of $f(x)$, then $\alpha^2 - 2$ is also a root of $f(x)$.}
        \textbf{Hint: }You might want to use Sage to help you with the algebra here.
        Use \texttt{R.<x> = QQ[]} to create a polynomial ring called \texttt{R} with a variable \texttt{x} and rational coefficients.
        You can use \texttt{f.subs(g)} to substitute the expression \texttt{g} for \texttt{x} in the polynomial \texttt{f}.
        The \texttt{quo\_rem} function will be useful.

        \item \blue{Prove that $f(x)$ must have at least one real root.}

        \item \blue{Prove that $f(x)$ is separable.}

        \item Let $\alpha \in \RR$ be such a root.
        \blue{Prove that $\QQ(\alpha)/\QQ$ is a splitting extension of the separable polynomial $f(x)$.}
        Hence $\QQ(\alpha)/\QQ$ is a Galois extension.

        \item \blue{What is the Galois group of $\QQ(\alpha)/\QQ$?
        Determine the structure of $\gal(\QQ(\alpha)/\QQ)$ and find explicit generators and relations.}
    \end{enumerate}
\end{problem}
%Examples


\end{document}